\documentclass[a4paper,10.5pt]{article}
\usepackage[T1]{fontenc}\usepackage[utf8]{inputenc}
\usepackage[top=2.2cm,bottom=2.2cm,left=2.35cm,right=2.35cm]{geometry}
\usepackage{hyperref}\usepackage{xcolor}\definecolor{accent}{HTML}{1F4E78}\hypersetup{colorlinks=true,urlcolor=accent}
\pagestyle{empty}\setlength{\parindent}{0pt}\setlength{\parskip}{7pt}
\begin{document}
\textbf{Prabhavit Gupta}\\
Venice, Italy \textbar{} +39 351 610 7807 \textbar{} prabhavitg@gmail.com\\[10pt]
Gentile team di Ticketoo,

mi candido per la posizione di Software Engineer a Roma. Ho tre anni di esperienza come sviluppatore full-stack e AI in otto aziende europee, sempre su sistemi in produzione con utenti reali. Il mio Python è soprattutto pipeline LLM, automazione, scraping e QA: non ho ancora usato Django in produzione, ma conosco FastAPI e imparo in fretta quando serve.

Per UltraStrategy ho costruito in Python la pipeline di parsing dei CV integrata nella piattaforma: l'utente carica il CV, la pipeline estrae i dati e ne valuta la rilevanza per il ruolo, il risultato arriva al frontend e infine genera un "Ultra CV" standardizzato, usando Affinda e le API di Claude.

Per Get2Germany ho integrato le Google API e costruito strumenti AI per trovare corsi, stage e offerte di lavoro, oltre a ottimizzare il database che li sosteneva.

Per FlowSpace ho automatizzato la migrazione dei dati da uno stack legacy verso un backend Laravel e Node, mappando la struttura vecchia sulla nuova ed eseguendo il cutover con il prodotto sempre in produzione, portando la piattaforma da 2.000 a oltre 4.500 utenti. Ho anche costruito integrazioni API con OfficeRnD e Nexudus per le prenotazioni automatizzate.

Quello che mi interessa del vostro annuncio è proprio questo equilibrio tra backend solido e sistemi anti-frode basati su AI, il tipo di lavoro end-to-end che ho già fatto. Vivo a Venezia, sono autorizzato a lavorare in Italia senza necessità di sponsorship, e sono disponibile a trasferirmi a Roma o a lavorare in modalità ibrida. Resto a vostra disposizione per la call conoscitiva.

Cordiali saluti,\\
Prabhavit Gupta
\end{document}
